\documentclass[UTF8,a4paper,12pt]{ctexart}

\usepackage{amsmath,amssymb,bm}
\usepackage{booktabs}
\usepackage{geometry}
\usepackage{longtable}
\usepackage{array}
\usepackage{enumitem}
\usepackage{microtype}
\usepackage{hyperref}

\geometry{left=2.6cm,right=2.6cm,top=2.5cm,bottom=2.5cm}
\setlength{\parindent}{2em}
\setlength{\parskip}{0.35em}
\renewcommand{\arraystretch}{1.35}
\setlist{nosep}
\hypersetup{hidelinks}

\title{方法与案例研究部分}
\author{}
\date{}

\begin{document}

\maketitle

\section{3 面向多线换乘站疏散的队列感知路径组织模型}

大型五线换乘站的应急疏散可以理解为多类乘客沿站内服务链向安全出口转移的过程。与单一站台或单一线路车站相比，换乘站内部同时存在站台出站流、列车到达释放流、站厅滞留流和跨线换乘流。不同来源客流并不是沿相互独立的路径离站，而是在楼梯、扶梯、换乘通道、闸机和出口前区等有限设施处反复交汇。因此，全站疏散效率不仅取决于几何路径长度，还取决于乘客批次真正到达共享设施时的排队等待和下游空间接收能力。

本文将换乘客流作为解释全站疏散过程的重要机制变量。换乘客流通常跨越多个线路空间，经过的服务链更长，且更容易与列车到达客流和普通出站客流在同一垂向设施、通道或闸机组处形成叠加。本文的研究目标不是为某一类乘客单独寻找最短换乘路径，而是在整体疏散框架下刻画换乘客流参与共享设施竞争后对尾部清空、队列扩散和出口利用的影响。

为此，本文构建队列感知的动态路径组织模型 AdaptiveQueueAwareAStar，简称 AA。模型由三部分组成：多线换乘站疏散服务链网络、共享设施到达时刻队列预测模型，以及车站疏散效果综合评价方法。前两部分回答“客流如何在站内网络中汇聚”和“乘客到达设施时会面对怎样的等待”，第三部分用于把路径组织结果转化为可解释的疏散效率、瓶颈暴露和跨模型对照指标。

\subsection{3.1 多线换乘站疏散服务链网络}

将车站公共区抽象为有向网络
\[
G=(V,E),
\]
其中，节点集合 \(V\) 表示站台、站厅、换乘节点、楼梯、扶梯、闸机、出口前区和地面出口等空间或设施单元，边集合 \(E\) 表示相邻单元之间允许乘客移动的方向性连接。每条边 \(e=(i,j)\) 记录长度 \(l_e\)、有效宽度 \(w_e\)、设施类型、行进方向和通行能力属性；每个节点 \(v\) 记录面积 \(A_v\)、当前人数 \(N_v(t)\)、节点类型和空间接收上限 \(\bar N_v\)。

根据乘客在站内的活动链条，客流来源组集合记为
\[
\mathcal{G}=\mathcal{G}^{out}\cup\mathcal{G}^{train}\cup\mathcal{G}^{hall}\cup\mathcal{G}^{tr},
\]
其中，\(\mathcal{G}^{out}\) 表示站台等待或出站客流，\(\mathcal{G}^{train}\) 表示列车到达释放客流，\(\mathcal{G}^{hall}\) 表示站厅滞留客流，\(\mathcal{G}^{tr}\) 表示换乘客流。对于换乘客流，模型保留来源线路、目标线路或最终离站目标，使其通过量、等待时间和出口使用能够与普通出站客流分开统计，但评价目标仍然是全站整体疏散。

在离散仿真中，一个进入路径组织阶段的客流批次记为
\[
b=(g,o_b,d_b,n_b,t_b),
\]
其中，\(g\in\mathcal{G}\) 为来源组，\(o_b\) 为批次当前节点，\(d_b\) 为可接受目标集合或目标方向，\(n_b\) 为批次人数，\(t_b\) 为批次进入路径组织阶段的时刻。采用批次建模可以保留列车到达和换乘释放的时间结构，也便于在路径搜索后登记其对共享设施的未来到达需求。

站内许多设施在图模型中可能被多条边引用，但在现实疏散中只具有一份物理通行能力。为避免同一楼梯、扶梯、闸机组或出口前区被重复释放容量，定义边到共享设施资源的映射
\[
\rho:E\rightarrow R\cup\{\varnothing\},
\]
其中，\(R\) 为有限通行设施资源集合。若 \(\rho(e)=r\)，则边 \(e\) 的通行需要消耗资源 \(r\) 的服务能力；若 \(\rho(e)=\varnothing\)，则该边仅表示普通空间移动。映射 \(\rho\) 使来自不同线路、不同站台和不同换乘关系的客流在同一资源上共同竞争容量，从而把换乘站的服务链交叉关系显式写入模型。

普通移动边的物理行走时间由边长和速度模型确定：
\[
\tau_e(t)=\frac{l_e}{v_e(t)}.
\]
其中，\(v_e(t)\) 可依据基础速度-密度关系、设施类型和仿真配置给出。速度-密度关系在本文中只作为基础运动模型的一部分，用于计算物理行走时间和拥挤暴露，不作为主要敏感性分析对象。

\subsection{3.2 共享设施到达时刻队列预测模型}

AA 的核心不是读取当前队列后立即选择一条最短路径，而是在候选路径评价时预测客流批次到达共享设施的时刻，并估计该时刻之前已经承诺到达该设施的其他客流需求。设共享设施 \(r\) 的服务率为 \(\mu_r\)，当前时刻为 \(t\)，设施入口等待队列为 \(Q_r(t)\)。在当前路径组织轮次中，已经登记但尚未完成服务的未来到达事件集合记为
\[
\mathcal{A}_r(t)=\{(\eta_j,n_j,g_j,b_j)\},
\]
其中，\(\eta_j\) 为事件 \(j\) 预计到达设施 \(r\) 的时刻，\(n_j\) 为到达人数，\(g_j\) 为来源组，\(b_j\) 为批次编号。这些事件来自两类对象：一类是在执行层中已经进入移动过程、预计稍后到达设施的批次；另一类是在同一路径组织轮次中，前序批次已经选择路径后登记的软性到达意图。

对候选批次 \(b\) 而言，若其沿候选路径预计在 \(\eta\) 时刻到达设施 \(r\)，则需要估计 \([t,\eta]\) 内队列的演化。将 \(\mathcal{A}_r(t)\) 中不晚于 \(\eta\) 的事件按时间排序，得到 \(s_1,\ldots,s_K\)，并记 \(s_0=t\)。若 \(a_k\) 为时刻 \(s_k\) 到达的总人数，则队列递推为
\[
q_k=\max\{q_{k-1}-\mu_r(s_k-s_{k-1}),0\}+a_k,\quad k=1,\ldots,K,
\]
其中，\(q_0=Q_r(t)\)。批次 \(b\) 到达设施 \(r\) 时的预测队列为
\[
\widehat Q_r(\eta)=\max\{q_K-\mu_r(\eta-s_K),0\}.
\]
若 \([t,\eta]\) 内没有已登记事件，则上式退化为当前队列在服务释放后的剩余量。由此得到批次到达设施时的预测等待时间
\[
W_r(\eta)=\frac{\widehat Q_r(\eta)}{\mu_r}.
\]

除前方队列等待外，批次自身也会占用设施服务能力。为了把批次人数对设施释放时间的影响纳入路径代价，设批次在资源 \(r\) 上形成的平均服务占用时间为
\[
S_r(n_b)=\frac{n_b}{2\mu_r}.
\]
该项表示一个批次通过有限服务设施时的平均排放时间贡献，用于避免多个大批次同时把同一设施视为可无成本通过。

对候选路径 \(P_b=(e_1,e_2,\ldots,e_m)\)，AA 按路径顺序递推到达时刻。设进入第 \(i\) 条边或其对应资源的时刻为 \(\eta_i\)，则路径目标函数写为
\[
J_b(P_b)=\sum_{i=1}^{m}\left[\tau_{e_i}(\eta_i)+I_{\rho(e_i)}\left(W_{\rho(e_i)}(\eta_i)+S_{\rho(e_i)}(n_b)\right)\right]+\lambda R_b(P_b),
\]
其中，\(I_{\rho(e_i)}\) 表示边 \(e_i\) 是否对应共享资源，\(R_b(P_b)\) 为路径上的拥挤暴露辅助项，\(\lambda\) 为其权重。该目标函数的主要作用是把物理行走时间、到达时刻设施队列、批次自身服务占用和拥挤暴露放在同一时间维度下评价。当前批次若选择某条路径，其预计到达共享设施的事件会写入 \(\mathcal{A}_r(t)\)，使后续批次能够感知前序批次对未来容量的占用。

\subsection{3.3 车站疏散效果综合评价方法}

本文评价体系分为整体效率、共享设施瓶颈、来源组机制和跨模型一致性四类。整体效率指标包括 \(T_{50}\)、\(T_{80}\)、\(T_{95}\)、\(T_{100}\)、平均疏散时间、累计停滞人秒、资源排队人秒、空间阻塞人秒、总移动距离和计算耗时。其中，\(T_{95}\) 和 \(T_{100}\) 用于刻画大部分乘客离站后的尾部清空阶段，累计停滞人秒用于衡量排队、阻塞和等待暴露。

共享设施瓶颈指标包括关键楼梯、扶梯、换乘通道、闸机组和出口前区的通过人数、峰值队列、队列人秒、首次排队时刻、最后排队时刻和设施利用率。该类指标用于判断路径组织是否只是改变了几何距离，还是改变了有限设施上的负荷分配和队列持续时间。

来源组机制指标按线路和客流类型统计。对每个来源组 \(g\)，记录其清空时间 \(T_g\)、关键设施通过量、等待暴露和出口使用分布。换乘客流不作为唯一研究对象，但其通过共享设施的路径更长、与其他来源交叉更多，因此在结果解释中单独报告换乘来源组的等待暴露和出口利用，有助于说明五线换乘结构如何影响全站疏散。

跨模型一致性指标用于比较图模型执行结果与 Pathfinder 微观仿真趋势。对照内容包括累计疏散曲线、主要拥堵位置、关键设施通行顺序和出口利用分布。由于图模型和 Pathfinder 的行人行为规则、空间离散方式和局部避让机制不同，本文将 Pathfinder 作为趋势对照和机制检验，而不把二者逐人逐秒完全等同。

\section{4 队列感知路径组织模型求解方法}

第 3 章给出了 AA 的模型构成，本章说明模型如何在离散疏散仿真中求解。求解过程的关键是把路径组织层和设施服务执行层放在同一时间推进框架内：路径组织层负责生成批次的路径意图，执行层负责在共享容量和下游接收约束下决定实际通行人数。二者分工不同，但必须使用同一套网络、同一套容量和同一套客流状态。

\subsection{4.1 共享容量约束下的动态加载转换}

设时刻 \(t\) 的系统状态为
\[
\Omega(t)=\left(G,\mathcal{B}(t),\mathcal{Q}(t),\mathcal{N}(t),\mathcal{A}(t)\right),
\]
其中，\(\mathcal{B}(t)\) 为当前可决策批次集合，\(\mathcal{Q}(t)\) 为共享设施队列状态，\(\mathcal{N}(t)\) 为节点和边上的客流占用状态，\(\mathcal{A}(t)\) 为已登记的未来到达事件。每个时间步内，模型先处理上一时间步已经完成的移动和到达事件，再刷新节点人数、边密度和资源队列，随后对可决策批次执行路径组织。

在同一时间步内，所有通过同一资源 \(r\) 的请求共同竞争该资源的服务能力。若仿真步长为 \(\Delta t\)，资源 \(r\) 在该时间步可释放的人数为
\[
C_r(t)=\left\lfloor \mu_r\Delta t+\xi_r(t)\right\rfloor ,
\]
其中，\(\xi_r(t)\) 为上一时间步保留的非整数服务能力余量。设本时间步请求进入资源 \(r\) 的人数为 \(A_r(t)\)，则队列更新为
\[
Q_r(t+\Delta t)=\max\{Q_r(t)+A_r(t)-C_r(t),0\}.
\]

动态加载还需要检查下游空间接收条件。设节点 \(v\) 的空间接收上限为 \(\bar N_v\)，当前占用为 \(N_v(t)\)，已经被其他移动预留的进入人数为 \(B_v(t)\)，则新进入人数 \(y_v(t)\) 应满足
\[
N_v(t)+B_v(t)+y_v(t)\leq \bar N_v .
\]
若下游空间不能接收，即使上游资源仍有名义服务能力，相关乘客也不能被释放到下游节点，而应保留在上游等待状态。这一处理用于刻画高负荷换乘站中常见的通道回溢和出口前区滞留。

共享容量动态加载的意义在于保证客流守恒和容量守恒。对同一时间步中的多来源请求，模型不会因为请求来自不同线路、不同路径或不同批次而重复释放同一设施容量；对未能获得服务能力或下游接收空间的乘客，模型将其保留在上游队列或阻塞状态，并在下一时间步继续参与竞争。

\subsection{4.2 AA 模型求解}

AA 对每个可决策批次执行一次时间依赖 A* 搜索。搜索标签记为
\[
L=(v,\eta,c,\pi),
\]
其中，\(v\) 为当前节点，\(\eta\) 为预计到达当前节点的时刻，\(c\) 为累计目标函数值，\(\pi\) 为前驱标签。与静态最短路不同，两个标签即使位于同一节点，只要到达时刻不同，也可能面对不同的未来设施队列，因此形式上不能简单合并为同一个静态节点状态。

标签从节点 \(u\) 沿边 \(e=(u,v)\) 扩展时，首先计算物理行走时间 \(\tau_e\)，随后判断该边是否对应共享资源。若 \(\rho(e)=r\)，则根据预计进入资源的时刻 \(\eta\) 查询 \(\widehat Q_r(\eta)\)，并计算预测等待 \(W_r(\eta)\) 和批次自身服务占用 \(S_r(n_b)\)。若下游空间接收条件可能限制移动，则加入相应的空间等待或阻塞代价。扩展完成后，标签的到达时刻和累计目标函数值同步更新。

AA 使用自由流条件下到最近可用出口的时间下界作为启发式函数，以减少搜索空间。启发式只提供剩余时间下界，不替代真实路径代价。真实路径代价仍由已递推的物理行走时间、到达时刻队列等待、批次服务占用和拥挤暴露共同确定。

批次路径被接受后，模型并不立即认为该批次已经通过相关设施，而是把路径上各共享资源的预计进入时刻、人数和来源组登记为未来到达事件。后续批次搜索候选路径时会读取这些事件，从而感知前序批次对同一设施未来服务能力的占用。该机制使路径组织不再基于同一份未更新状态独立决策，而是在一个可复现的顺序中逐步形成共享设施负荷预期。

对已经进入移动过程、已经获得设施服务能力或处于不可重评估阶段的批次，AA 继续执行既有移动，不在途中任意改变方向。对仍处于路径选择阶段的批次，若候选新路径相对于保留路径具有足够收益，且满足路径有效性和方向稳定条件，则允许重新评价路径。设保留路径剩余代价为 \(C_{\mathrm{stay}}\)，候选路径代价为 \(C_{\mathrm{switch}}\)，相对收益为
\[
R=\frac{C_{\mathrm{stay}}-C_{\mathrm{switch}}}{C_{\mathrm{stay}}}.
\]
只有当 \(C_{\mathrm{switch}}<C_{\mathrm{stay}}\) 且 \(R\) 超过内部防摆动阈值时，路径切换才被接受。该阈值属于算法稳定性参数，不作为行人行为参数解释。

本文使用 ImprovedAStar 作为计算对照方法。两种方法采用相同的站内网络、相同的客流输入、相同的速度模型、相同的共享设施容量、相同的下游接收约束和相同的终止条件。二者差异限定在路径组织层：AA 显式登记未来到达事件并预测到达时刻队列，ImprovedAStar 使用当前网络状态和即时边权进行路径评价，不引入已承诺未来到达负荷。ImprovedAStar 仅为计算基准，不对应龙阳路站已发布或实际执行的现场疏散安排。

\section{5 案例研究}

本文以龙阳路站作为案例，检验 AA 在大型多线换乘站整体应急疏散中的作用。龙阳路站连接 2 号线、7 号线、16 号线、18 号线和磁浮线，站内包含多线路站台、站厅、换乘通道、楼扶梯、闸机和多个地面出口。五线换乘结构使不同来源客流在站内服务链上具有明显交叉：站台出站客流、列车到达释放客流和换乘客流会在部分垂向设施、换乘连接和出站设施处共同加载。

案例模型依据 CAD 图纸和仿真配置建立站内拓扑。几何对象包括站台、站厅、换乘节点、楼扶梯、闸机、出口前区和地面出口；客流对象包括各线路站台等待客流、列车到达释放客流、站厅滞留客流和换乘关系客流。未经过核验的设施容量不写作官方数据；最终参数表需将 CAD/模型配置值、文献取值和必要假设值分开标注。

实验设置低负荷常规应急场景和高负荷满载列车叠加场景。低负荷常规应急场景包含站台等待、站厅滞留和基础换乘客流，输入总人数为 2187 人，用于检验共享设施竞争较弱时 AA 是否保持稳定、是否引入不必要绕行或出口偏置。高负荷满载列车叠加场景在基础客流上叠加多线路列车到达释放，输入总人数为 17905 人，用于放大多来源客流在共享设施上的同步汇聚效应，并检验到达时刻队列预测对尾部清空和关键设施排队的影响。

\begin{table}[htbp]
\centering
\caption{低负荷与高负荷疏散场景设置}
\small
\begin{tabular}{p{0.20\textwidth}p{0.14\textwidth}p{0.28\textwidth}p{0.28\textwidth}}
\toprule
场景 & 输入人数 & 主要客流构成 & 实验目的 \\
\midrule
低负荷常规应急 & 2187 & 站台等待、站厅滞留和基础换乘客流 & 检验设施竞争较弱时 AA 的稳定性和低负荷适用性 \\
高负荷满载列车叠加 & 17905 & 基础客流叠加多线路列车到达释放客流 & 检验多来源客流同步汇聚时 AA 对队列等待和尾部清空的影响 \\
\bottomrule
\end{tabular}
\end{table}

由于最终仿真结果尚未冻结，以下结果表保留指标结构和解释口径。正式数值应在低负荷、高负荷、敏感性分析和 Pathfinder 对照结果统一冻结后填入。

\subsection{5.1 到达负荷与队列预测结果分析}

案例研究首先分析 AA 的中间预测结果，而不是直接进入总疏散时间比较。该安排用于验证模型核心机制：AA 是否能够在路径选择阶段识别乘客批次到达关键共享设施时的未来排队等待。若该中间机制不能被说明，后续关于整体疏散效率的比较就缺乏可解释性。

本节选取换乘通道、楼扶梯、闸机组和出口前区中的关键共享设施，比较预测到达人数、预测队列、执行层实际队列、实际通过人数和未释放人数。对每个关键设施，记录队列峰值、峰值出现时刻、首次排队时刻、最后排队时刻和队列人秒。若预测队列与执行层实际队列在峰值时段和瓶颈位置上具有一致趋势，说明到达时刻负荷预测能够捕捉共享设施竞争；若差异较大，则需要从批次粒度、服务率取整、下游接收限制或微观行为差异解释。

\begin{table}[htbp]
\centering
\caption{到达负荷与队列预测结果记录表}
\small
\begin{tabular}{p{0.18\textwidth}p{0.13\textwidth}p{0.13\textwidth}p{0.13\textwidth}p{0.13\textwidth}p{0.20\textwidth}}
\toprule
设施类型 & 预测峰值队列 & 实际峰值队列 & 首次排队时刻 & 最后排队时刻 & 主要来源组构成 \\
\midrule
关键楼扶梯 & 待冻结 & 待冻结 & 待冻结 & 待冻结 & 待冻结 \\
换乘通道 & 待冻结 & 待冻结 & 待冻结 & 待冻结 & 待冻结 \\
闸机组 & 待冻结 & 待冻结 & 待冻结 & 待冻结 & 待冻结 \\
出口前区 & 待冻结 & 待冻结 & 待冻结 & 待冻结 & 待冻结 \\
\bottomrule
\end{tabular}
\end{table}

对代表性批次，进一步报告候选路径的代价分解。若 AA 改变某一批次路径，应说明变化来自物理行走时间、设施等待、批次服务占用还是下游空间接收限制。该分析用于区分“绕行换取排队减少”和“几何距离缩短”两类不同机制，避免把最终总时间变化简单归因于算法名称。

\subsection{5.2 路径组织方案性能分析}

在到达负荷与队列预测结果得到说明后，本节比较 AA 与 ImprovedAStar 的路径组织方案性能。比较采用相同站内网络、相同客流输入和相同执行层规则，因此结果差异主要来自路径组织逻辑。低负荷场景重点观察 AA 是否保持稳定，避免在共享设施尚未形成明显竞争时造成额外绕行；高负荷场景重点观察 AA 是否通过提前感知未来队列改变关键设施负荷分布，从而影响尾部清空和停滞等待。

\begin{table}[htbp]
\centering
\caption{低负荷与高负荷场景下整体疏散性能记录表}
\small
\begin{tabular}{p{0.16\textwidth}p{0.13\textwidth}p{0.09\textwidth}p{0.10\textwidth}p{0.10\textwidth}p{0.12\textwidth}p{0.12\textwidth}p{0.10\textwidth}}
\toprule
场景 & 方法 & 人数 & \(T_{95}\) & \(T_{100}\) & 累计停滞人秒 & 总移动距离 & 计算耗时 \\
\midrule
低负荷常规应急 & ImprovedAStar & 2187 & 待冻结 & 待冻结 & 待冻结 & 待冻结 & 待冻结 \\
低负荷常规应急 & AA & 2187 & 待冻结 & 待冻结 & 待冻结 & 待冻结 & 待冻结 \\
高负荷满载列车叠加 & ImprovedAStar & 17905 & 待冻结 & 待冻结 & 待冻结 & 待冻结 & 待冻结 \\
高负荷满载列车叠加 & AA & 17905 & 待冻结 & 待冻结 & 待冻结 & 待冻结 & 待冻结 \\
\bottomrule
\end{tabular}
\end{table}

整体指标之后，需要进一步分析路径组织方案在空间上的变化。主要比较内容包括各出口利用人数、各线路或来源组清空时间、关键设施通过量、资源排队人秒和空间阻塞人秒。若 AA 缩短尾部清空时间，应进一步说明这种变化是否来自关键楼扶梯、换乘通道、闸机或出口前区的持续排队减少；若 AA 降低队列等待但增加移动距离或计算耗时，也应同时报告该权衡。

本节不使用单一综合分数替代具体指标。原因是全站疏散效果具有多维性：总疏散时间、尾部清空、关键设施队列、出口均衡和计算代价可能并不同时改善。正式结果写作时，应分别说明效率改善、拥堵缓解、出口利用变化和计算开销，避免把局部指标改进泛化为整体最优。

\subsection{5.3 路径组织与乘客疏散效果集成分析}

本节将路径组织结果放回完整疏散执行过程，分析 AA 是否改变了多来源客流在站内服务链上的加载方式。分析重点不只是总人数何时清空，而是客流如何经过关键共享设施、哪些来源组参与瓶颈形成，以及图模型结论是否能在 Pathfinder 微观仿真中得到趋势支持。

\subsubsection{5.3.1 关键设施与换乘客流机制分解}

对关键设施，报告通过人数、峰值队列、队列人秒、首次排队和最后排队时刻。对来源组，分别统计站台出站、列车到达、站厅滞留和换乘客流在关键设施处的通过量和等待暴露。换乘客流的分析重点不是证明换乘客流自身更快离站，而是解释其跨线服务链如何与其他来源客流共同影响全站队列分布。

\begin{table}[htbp]
\centering
\caption{关键设施和来源组机制分解记录表}
\small
\begin{tabular}{p{0.18\textwidth}p{0.16\textwidth}p{0.13\textwidth}p{0.13\textwidth}p{0.13\textwidth}p{0.17\textwidth}}
\toprule
分析对象 & 方法 & 通过人数 & 峰值队列 & 队列人秒 & 主要来源组 \\
\midrule
关键楼扶梯 & ImprovedAStar & 待冻结 & 待冻结 & 待冻结 & 待冻结 \\
关键楼扶梯 & AA & 待冻结 & 待冻结 & 待冻结 & 待冻结 \\
换乘通道 & ImprovedAStar & 待冻结 & 待冻结 & 待冻结 & 待冻结 \\
换乘通道 & AA & 待冻结 & 待冻结 & 待冻结 & 待冻结 \\
闸机组 & ImprovedAStar & 待冻结 & 待冻结 & 待冻结 & 待冻结 \\
闸机组 & AA & 待冻结 & 待冻结 & 待冻结 & 待冻结 \\
\bottomrule
\end{tabular}
\end{table}

若结果显示 AA 将部分客流从某一持续排队设施转移至服务能力尚有余量的路径，应结合路径代价分解说明其原因：前者可能几何距离更短，但到达时刻队列更长；后者可能步行距离略长，但等待时间和尾部滞留更低。若结果显示低负荷场景下两种方法差异较小，也不应强行解释为缺陷，而应说明在共享设施竞争不明显时，到达时刻队列预测的边际作用有限。

\subsubsection{5.3.2 Pathfinder 微观仿真对照}

Pathfinder 对照用于检验图模型路径组织结果在微观行人仿真中的趋势一致性。对照不改变本文提出的 AA 模型，也不把 Pathfinder 作为 AA 的求解器；其作用是在不同建模粒度下比较疏散曲线、拥堵位置和出口利用是否呈现相近趋势。

Pathfinder 场景使用与图模型一致的站内空间、客流规模和路径组织输出。对每种方法，按来源组或路径组生成对应的行人行为与控制路径，并在微观仿真中记录总疏散时间、个人离站时间分布、拥堵持续时间、楼扶梯和闸机周边聚集情况以及出口使用分布。若 Pathfinder 中的拥堵位置与图模型识别的关键设施一致，且累计疏散曲线和尾部阶段趋势相近，则可作为 AA 机制有效性的独立支撑；若两者存在偏差，应从微观避让、局部几何细节、行人异质性和路径执行差异解释。

\begin{table}[htbp]
\centering
\caption{Pathfinder 对照分析指标}
\small
\begin{tabular}{p{0.22\textwidth}p{0.34\textwidth}p{0.32\textwidth}}
\toprule
对照维度 & 图模型输出 & Pathfinder 输出 \\
\midrule
疏散曲线 & 累计疏散人数、\(T_{95}\)、\(T_{100}\) & 个人离站时间分布和累计疏散曲线 \\
拥堵位置 & 关键设施队列峰值和队列人秒 & 密度云图、拥堵持续时间和局部聚集区域 \\
出口利用 & 各出口通过人数及来源组构成 & 各出口实际离站人数和占比 \\
来源组差异 & 换乘、站台、列车到达和站厅客流清空时间 & 不同人群组离站时间和拥堵暴露 \\
\bottomrule
\end{tabular}
\end{table}

\subsubsection{5.3.3 适用边界与敏感性分析}

敏感性分析服务于本文机制，而不是泛泛测试所有行人参数。本文重点分析三类因素：换乘客流比例、列车到达重叠程度和关键共享设施服务能力。换乘客流比例用于检验服务链交叉强度变化时 AA 的作用边界；列车到达重叠程度用于检验多来源客流在时间上集中释放时的队列预测价值；关键共享设施服务能力用于检验瓶颈强弱变化对路径组织效果的影响。

\begin{table}[htbp]
\centering
\caption{敏感性分析设置与解释目的}
\small
\begin{tabular}{p{0.24\textwidth}p{0.30\textwidth}p{0.34\textwidth}}
\toprule
敏感性因素 & 调整方式 & 解释目的 \\
\midrule
换乘客流比例 & 在总需求或基础需求不变条件下调整换乘来源组占比 & 检验服务链交叉增强时 AA 是否更能识别共享设施未来排队 \\
列车到达重叠程度 & 调整多线路列车释放的时间间隔或同步程度 & 检验客流到达高峰越集中时到达时刻预测是否更重要 \\
关键共享设施服务能力 & 调整代表性楼扶梯、通道、闸机或出口前区容量 & 检验设施瓶颈强弱变化下 AA 的适用范围和失效边界 \\
\bottomrule
\end{tabular}
\end{table}

最终结果写作时，若 AA 的优势主要出现在高换乘比例、高到达重叠或关键设施服务能力较低的条件下，应将其解释为面向共享设施竞争明显场景的路径组织方法；若在低负荷或瓶颈较弱场景下差异较小，则说明模型不会在非瓶颈条件下强行产生显著改善。该解释比单纯报告高低负荷结果更能说明 AA 适用于哪些多线换乘站疏散情境。

\end{document}
